\PilotPage{7}
\section*{From complex exponentials to trigonometric functions}

\begin{equation}
    C_1e^{i\alpha}+C_2e^{-i\alpha}=A\cos\alpha+B\sin\alpha.
    \tag{12a}
\end{equation}
or
\begin{equation}
    C_1e^{i\alpha}+C_2e^{-i\alpha}=C\sin(\alpha+\varphi).
    \tag{12b}
\end{equation}

\subsection*{Proof}

Recall Euler's identity,
\begin{equation*}
    e^{i\alpha}=\cos\alpha+i\sin\alpha,
\end{equation*}
and hence
\begin{equation*}
    e^{-i\alpha}=\cos\alpha-i\sin\alpha.
\end{equation*}
Therefore,
\begin{align*}
    C_1e^{i\alpha}&=C_1\cos\alpha+iC_1\sin\alpha,\\
    C_2e^{-i\alpha}&=C_2\cos\alpha-iC_2\sin\alpha.
\end{align*}
Adding,
\begin{align*}
    C_1e^{i\alpha}+C_2e^{-i\alpha}
      &=(C_1+C_2)\cos\alpha+i(C_1-C_2)\sin\alpha\\
      &=A\cos\alpha+B\sin\alpha,
\end{align*}
where
\begin{equation*}
    A=C_1+C_2,\qquad B=i(C_1-C_2).
\end{equation*}

Let
\begin{equation*}
    C=\sqrt{A^2+B^2},\qquad \varphi=\tan^{-1}\frac{A}{B}.
\end{equation*}
Recall that
\begin{equation*}
    \sin(\alpha+\varphi)=\sin\varphi\cos\alpha+\cos\varphi\sin\alpha.
\end{equation*}
Hence
\begin{equation*}
    C\sin(\alpha+\varphi)=(C\sin\varphi)\cos\alpha+(C\cos\varphi)\sin\alpha.
\end{equation*}
Since
\begin{equation*}
    A^2+B^2=C^2\sin^2\varphi+C^2\cos^2\varphi=C^2,
\end{equation*}
we have $C=\sqrt{A^2+B^2}$, and
\begin{equation*}
    \frac{A}{B}=\frac{C\sin\varphi}{C\cos\varphi}=\tan\varphi
    \quad\Longrightarrow\quad
    \varphi=\tan^{-1}\frac{A}{B}.
\end{equation*}
Thus,
\begin{equation}
    C_1e^{i\alpha}+C_2e^{-i\alpha}=C\sin(\alpha+\varphi).
    \tag{12c}
\end{equation}
\hfill Q.E.D.
